% ============================================================
% Methodology section for:
% "A From-Scratch Swin Transformer Encoder--Decoder Network for
%  Automated Radiology Image Captioning"
%
% Paste this \section directly after Related Work
% (\label{sec:method}, as referenced in the Introduction's
% roadmap paragraph). Content is drawn directly from the actual
% implementation: swin_model.py, decoder_model.py,
% caption_model.py, vocab.py, and train.py.
% ============================================================

\section{Methodology}
\label{sec:method}

The proposed system follows an encoder--decoder design. A Swin Transformer encoder maps a $224\times224\times3$ radiology image to a fixed-size set of visual tokens, and an autoregressive Transformer decoder consumes those tokens as memory to generate a caption one word at a time. Every component -- patch embedding, windowed and shifted-window attention, patch merging, masked self-attention, cross-attention, and the training loop itself -- is implemented directly with primitive \texttt{torch.nn} layers, with no pretrained weights and no third-party Transformer or vision library.

\subsection{Swin Transformer Encoder}
\label{sec:method-encoder}

\subsubsection{Patch Embedding}
The input image is first divided into non-overlapping $4\times4$ patches using a stride-4 convolution (\texttt{PatchEmbed}), which is mathematically equivalent to flattening each patch and applying a shared linear projection. This produces a $56\times56$ grid of tokens with embedding dimension $C=96$, followed by layer normalization.

\subsubsection{Window and Shifted-Window Self-Attention}
Rather than computing self-attention across all $56\times56=3136$ tokens -- which would scale quadratically with token count -- the encoder partitions the token grid into non-overlapping $7\times7$ windows and restricts multi-head self-attention to within each window (\texttt{WindowAttention}). A learned relative position bias is added to the attention logits: for every possible relative offset $(\Delta h, \Delta w)$ between two positions inside a window, a bias term is looked up from a table of shape $(2\cdot7-1)^2 \times \text{heads}$ and added to the corresponding attention score, replacing the fixed absolute positional encoding used in the original Transformer.

Restricting attention to fixed windows would prevent tokens from ever exchanging information across window boundaries. Following the shifted-window design, alternating blocks within each stage cyclically shift the token grid by $\lfloor \text{window\_size}/2 \rfloor$ pixels before partitioning (\texttt{SwinBlock}, even-indexed blocks use regular W-MSA with \texttt{shift\_size=0}; odd-indexed blocks use shifted SW-MSA). Because the shift can create windows that mix content from non-adjacent regions of the original grid, an attention mask is precomputed per shifted block (\texttt{\_build\_shift\_mask}) that assigns a large negative bias ($-100.0$, rather than $-\infty$, to avoid a known softmax NaN issue on the Apple MPS backend) to attention scores between tokens that were not spatially adjacent before the shift. Each block applies pre-norm residual connections around both the (shifted-)window attention and a two-layer GELU feed-forward MLP.

\subsubsection{Patch Merging and Hierarchical Stages}
After every stage except the last, a patch-merging layer (\texttt{PatchMerging}) concatenates each $2\times2$ neighborhood of tokens along the channel dimension, applies layer normalization, and linearly projects back down to $2C$ channels -- halving spatial resolution while doubling channel width. The encoder stacks four such stages (\texttt{BasicLayer}) with depths $(2,2,6,2)$ and attention head counts $(3,6,12,24)$, following the Swin-T configuration. Starting from $56\times56\times96$, the four stages progressively reduce the representation to $28\times28\times192 \rightarrow 14\times14\times384 \rightarrow 7\times7\times768$. A final layer normalization is applied, yielding $49$ visual tokens of dimension $768$, which serve as the memory input to the decoder.

\subsection{Transformer Decoder}
\label{sec:method-decoder}

The decoder (\texttt{CaptionDecoder}) is a six-layer autoregressive Transformer that generates the caption token by token, conditioned on the 49 Swin visual tokens.

\subsubsection{Token Embedding and Positional Encoding}
Caption tokens are first mapped to $d_{model}=768$-dimensional vectors through a learned embedding table, matching the encoder's output dimension so no additional projection is needed between encoder and decoder. Fixed sinusoidal positional encodings (\texttt{PositionalEncoding}) are added to the token embeddings to inject sequence-order information, since self-attention itself is permutation-invariant.

\subsubsection{Decoder Layer: Masked Self-Attention, Cross-Attention, Feed-Forward}
Each of the six decoder layers (\texttt{DecoderLayer}) applies three sub-layers in sequence, each wrapped in a residual connection followed by post-layer-norm:
\begin{enumerate}
    \item \textbf{Masked self-attention} over the caption generated so far, using a lower-triangular causal mask so that position $t$ can only attend to positions $\le t$;
    \item \textbf{Cross-attention}, in which caption token representations act as queries while the 49 Swin encoder tokens act as keys and values -- this is the mechanism through which visual information enters the language stream, allowing each generated word to attend to the relevant regions of the image;
    \item \textbf{Position-wise feed-forward network}, a two-layer GELU MLP with hidden dimension $4\times d_{model}=3072$.
\end{enumerate}
All three attention operations (self-attention, cross-attention, and the encoder's window attention) share the same manually implemented multi-head attention routine, which scales dot-product scores by $1/\sqrt{d_k}$ and, when a mask is supplied, fills masked positions with $-100.0$ before the softmax -- the same numerically stable masking convention used throughout the encoder, applied consistently across both networks. The final decoder output is projected through a linear layer to vocabulary-sized logits at every position.

\subsection{Tokenization and Vocabulary}
\label{sec:method-vocab}

Captions are tokenized with a lightweight regex-based tokenizer (\texttt{Vocab.tokenize}): text is lowercased, all characters outside \texttt{[a-z0-9]} are stripped, and the result is whitespace-split. No subword tokenization (e.g., BPE) is used. A vocabulary is built from the training captions only, keeping every token with frequency $\ge 2$ and adding four special tokens (\texttt{<pad>}, \texttt{<sos>}, \texttt{<eos>}, \texttt{<unk>}); the final vocabulary is sorted alphabetically for a stable index mapping. Each caption is encoded as \texttt{<sos>} + tokens (truncated to \texttt{max\_len}$-2$) + \texttt{<eos>}, right-padded to a fixed length of 40 tokens.

\subsection{Training Procedure}
\label{sec:method-training}

The full model (\texttt{SwinCaptioningModel}) is trained end-to-end with teacher forcing: at each training step, the decoder is given the ground-truth caption shifted right as input and is trained to predict the ground-truth caption shifted left as the target, so that every position is supervised by the actual next word rather than the model's own previous prediction. The training objective is cross-entropy with label smoothing ($\varepsilon=0.1$), computed over all non-padding positions (\texttt{ignore\_index} set to the \texttt{<pad>} id).

Optimization uses AdamW with a peak learning rate of $5\times10^{-4}$, linear warmup over the first 5 epochs followed by cosine decay to zero over the remaining epochs, and gradient-norm clipping at $1.0$. Following standard practice for Transformer and Swin-style models, weight decay ($0.05$) is applied only to $\ge 2$-dimensional weight matrices; all 1-D parameters (LayerNorm weights/biases, linear biases) and the encoder's relative-position-bias table are placed in a separate parameter group with weight decay disabled, since decaying these distorts normalization scale and the learned window position bias without improving generalization. Training uses a batch size of 32, runs for up to 60 epochs, and stops early if validation loss does not improve for 6 consecutive epochs; the checkpoint with the best validation loss is retained for evaluation and inference. Input images are resized to $224\times224$, augmented during training with small random rotations ($\pm5^{\circ}$) and color jitter, and normalized with standard ImageNet channel statistics.

\subsection{Inference: Autoregressive Generation}
\label{sec:method-inference}

At inference time, the model generates captions autoregressively (\texttt{SwinCaptioningModel.generate}): the encoder computes the 49 visual tokens once, and the decoder is repeatedly queried starting from a single \texttt{<sos>} token, appending the greedily chosen highest-probability next token at each step (\texttt{argmax} decoding) until an \texttt{<eos>} token is produced or a maximum length of 40 tokens is reached. Two decoding-time constraints are applied to the logits before each \texttt{argmax}: the immediately preceding token is masked out to prevent trivial immediate repetition, and any token that would complete a trigram already seen earlier in the sequence is masked out (trigram blocking), reducing repeated phrases in the generated caption without requiring beam search.
